\documentclass[12pt,a4paper]{article}
\usepackage[utf8]{inputenc}
\usepackage[vietnamese]{babel}
\usepackage{amsmath,amssymb}
\usepackage[margin=2cm]{geometry}
\usepackage{enumitem}
\usepackage{fontspec}
\usepackage{booktabs}
\setmainfont{Times New Roman}

\title{\textbf{BÀI TẬP VỀ NHÀ: BÀI 7 -- NHIỆT HÓA HƠI RIÊNG}\\
\large Vật Lý 12 -- Chương trình GEMS}
\author{Giáo viên: Kha Khung Hiệp}
\date{}

\begin{document}
\maketitle
\textit{Thời gian làm bài: 45 phút}

\vspace{1em}
\hrule
\vspace{1em}

\section*{CÁC HẰNG SỐ CẦN NHỚ}
\begin{center}
\begin{tabular}{lccc}
\toprule
\textbf{Chất} & \textbf{Nhiệt hóa hơi riêng $L$ (kJ/kg)} & \textbf{Nhiệt dung riêng $c$ (kJ/(kg$\cdot$K))} & \textbf{Nhiệt độ sôi ($^{\circ}$C)} \\
\midrule
Nước & 2260 & 4,2 & 100 \\
Cồn (Ethanol) & 860 & 2,5 & 78 \\
Nhôm & -- & 0,9 & -- \\
Đồng & -- & 0,39 & -- \\
\bottomrule
\end{tabular}
\end{center}
\textit{Lưu ý: $1\text{ kJ} = 1000\text{ J}$; $1\text{ MJ} = 10^6\text{ J}$.}

\section*{PHẦN I: CÂU HỎI TRẮC NGHIỆM NHIỀU LỰA CHỌN (18 CÂU)}

\subsection*{A. Nhận biết (6 câu)}

\textbf{Câu 1:} Phát biểu nào sau đây là định nghĩa đúng về nhiệt hóa hơi riêng của một chất lỏng?
\begin{enumerate}[label=\Alph*.]
  \item Là nhiệt lượng cần cung cấp để 1 kg chất lỏng đó hóa hơi hoàn toàn ở nhiệt độ xác định.
  \item Là nhiệt lượng cần cung cấp để làm nóng 1 kg chất lỏng đó tăng thêm 1°C.
  \item Là nhiệt lượng tỏa ra khi 1 kg hơi của chất đó ngưng tụ hoàn toàn ở nhiệt độ xác định.
  \item Là năng lượng liên kết giữa các phân tử chất lỏng ở trạng thái sôi.
\end{enumerate}

\textbf{Câu 2:} Đơn vị đo của nhiệt hóa hơi riêng trong hệ SI là:
\begin{enumerate}[label=\Alph*.]
  \item $\text{J}/(\text{kg}\cdot\text{K})$.
  \item $\text{J}/\text{kg}$.
  \item $\text{J}$.
  \item $\text{W}/\text{kg}$.
\end{enumerate}

\textbf{Câu 3:} Nhiệt lượng cần cung cấp để một khối lượng $m$ chất lỏng hóa hơi hoàn toàn ở nhiệt độ sôi được xác định bằng công thức:
\begin{enumerate}[label=\Alph*.]
  \item $Q = L \cdot m$.
  \item $Q = m \cdot c \cdot \Delta t$.
  \item $Q = L \cdot m \cdot \Delta t$.
  \item $Q = \dfrac{L}{m}$.
\end{enumerate}

\textbf{Câu 4:} Ở áp suất tiêu chuẩn, nước sôi và hóa hơi ở nhiệt độ:
\begin{enumerate}[label=\Alph*.]
  \item $0^{\circ}\text{C}$.
  \item $100^{\circ}\text{C}$.
  \item $273\text{ K}$.
  \item $373^{\circ}\text{C}$.
\end{enumerate}

\textbf{Câu 5:} Nhiệt hóa hơi riêng của một chất lỏng đặc trưng cho:
\begin{enumerate}[label=\Alph*.]
  \item Thể tích và áp suất ngoài của chất lỏng.
  \item Bản chất của chất lỏng đó ở một nhiệt độ xác định.
  \item Khối lượng và nhiệt độ ban đầu của khối chất lỏng.
  \item Khả năng truyền nhiệt nhanh hay chậm của chất lỏng.
\end{enumerate}

\textbf{Câu 6:} Khi nấu cơm, ta thấy hơi nước bốc lên gặp nắp vung lạnh tạo thành các giọt nước bám ở mặt trong của nắp. Hiện tượng này được gọi là:
\begin{enumerate}[label=\Alph*.]
  \item Sự hóa hơi.
  \item Sự nóng chảy.
  \item Sự đông đặc.
  \item Sự ngưng tụ.
\end{enumerate}

\subsection*{B. Thông hiểu (6 câu)}

\textbf{Câu 7:} Phát biểu nào sau đây là đúng về mối quan hệ giữa nhiệt hóa hơi riêng và nhiệt độ?
\begin{enumerate}[label=\Alph*.]
  \item Nhiệt hóa hơi riêng không phụ thuộc vào nhiệt độ.
  \item Nhiệt hóa hơi riêng của một chất tăng khi nhiệt độ tăng.
  \item Nhiệt hóa hơi riêng của một chất tăng khi nhiệt độ giảm.
  \item Nhiệt hóa hơi riêng luôn bằng hằng số ở mọi nhiệt độ.
\end{enumerate}

\textbf{Câu 8:} Thiết bị nào sau đây hoạt động dựa trên nguyên lý thu nhiệt khi hóa hơi của chất lỏng?
\begin{enumerate}[label=\Alph*.]
  \item Bếp điện đun nước sôi.
  \item Máy điều hòa nhiệt độ (dàn bay hơi).
  \item Nồi hấp tiệt trùng y tế.
  \item Bình giữ nhiệt.
\end{enumerate}

\textbf{Câu 9:} Trong thí nghiệm đo nhiệt hóa hơi riêng của nước, người ta sử dụng oát kế để đo trực tiếp đại lượng nào sau đây?
\begin{enumerate}[label=\Alph*.]
  \item Khối lượng nước đã bay hơi.
  \item Nhiệt độ sôi của nước.
  \item Công suất dòng điện và thời gian đun nước.
  \item Nhiệt hóa hơi riêng của nước.
\end{enumerate}

\textbf{Câu 10:} Để hạn chế sai số do tỏa nhiệt ra môi trường trong thí nghiệm đo $L$ của nước, nhóm học sinh cần thực hiện biện pháp nào?
\begin{enumerate}[label=\Alph*.]
  \item Đun nước trong cốc thủy tinh không nắp.
  \item Dùng bình nhiệt lượng kế nhựa có vỏ xốp cách nhiệt tốt và nắp đậy kín.
  \item Sử dụng bếp ga đun nước thay cho dây đun điện trở.
  \item Đổ nước lạnh vào bình để thời gian đun sôi lâu hơn.
\end{enumerate}

\textbf{Câu 11:} Ý nghĩa vật lý của hằng số nhiệt hóa hơi riêng của nước $L = 2{,}26 \cdot 10^6\text{ J/kg}$ là gì?
\begin{enumerate}[label=\Alph*.]
  \item Cần cung cấp nhiệt lượng $2{,}26 \cdot 10^6\text{ J}$ để đun sôi 1 kg nước.
  \item Mỗi kilôgam nước cần thu nhiệt lượng $2{,}26 \cdot 10^6\text{ J}$ để hóa hơi hoàn toàn ở nhiệt độ sôi và áp suất chuẩn.
  \item Mỗi kilôgam nước tỏa ra nhiệt lượng $2{,}26 \cdot 10^6\text{ J}$ khi hóa hơi.
  \item Cần công suất $2{,}26 \cdot 10^6\text{ W}$ để hóa hơi hoàn toàn 1 kg nước sôi.
\end{enumerate}

\textbf{Câu 12:} Một học sinh đổ cồn (ethanol) lên da tay và đứng trước quạt thấy da mát buốt. Nguyên nhân là vì:
\begin{enumerate}[label=\Alph*.]
  \item Cồn bay hơi nhanh, thu nhiệt lượng lớn từ da tay để thực hiện quá trình hóa hơi.
  \item Cồn có nhiệt độ ban đầu rất thấp khi tiếp xúc với không khí.
  \item Cồn phản ứng hóa học tỏa nhiệt lạnh lên da tay.
  \item Gió quạt làm cho cồn ngưng tụ lại trên da.
\end{enumerate}

\subsection*{C. Vận dụng (6 câu)}

\textbf{Câu 13:} Nhiệt lượng cần thiết để hóa hơi hoàn toàn $0{,}5\text{ kg}$ nước ở nhiệt độ sôi $100^{\circ}\text{C}$ là bao nhiêu?
\begin{enumerate}[label=\Alph*.]
  \item $1130\text{ kJ}$.
  \item $2260\text{ kJ}$.
  \item $4520\text{ kJ}$.
  \item $565\text{ kJ}$.
\end{enumerate}

\textbf{Câu 14:} Cần cung cấp một nhiệt lượng bao nhiêu để hóa hơi hoàn toàn $3\text{ kg}$ cồn ở nhiệt độ sôi (nhiệt hóa hơi riêng của cồn $L = 8{,}6 \cdot 10^5\text{ J/kg}$)?
\begin{enumerate}[label=\Alph*.]
  \item $2{,}58 \cdot 10^6\text{ J}$.
  \item $8{,}60 \cdot 10^5\text{ J}$.
  \item $2{,}87 \cdot 10^5\text{ J}$.
  \item $5{,}16 \cdot 10^6\text{ J}$.
\end{enumerate}

\textbf{Câu 15:} Người ta đun sôi và hóa hơi hoàn toàn $2\text{ kg}$ nước ở $20^{\circ}\text{C}$. Nhiệt lượng tổng cộng cần cung cấp là bao nhiêu? (Biết $c = 4{,}2\text{ kJ/(kg}\cdot\text{K)}$, $L = 2260\text{ kJ/kg}$).
\begin{enumerate}[label=\Alph*.]
  \item $4520\text{ kJ}$.
  \item $672\text{ kJ}$.
  \item $5192\text{ kJ}$.
  \item $3848\text{ kJ}$.
\end{enumerate}

\textbf{Câu 16:} Để xác định gần đúng nhiệt hóa hơi riêng của nước, học sinh đun $0{,}2\text{ kg}$ nước ở $20^{\circ}\text{C}$ bằng bếp điện. Thời gian đun nước sôi (đến $100^{\circ}\text{C}$) là $4\text{ phút}$. Thời gian đun sôi để có $12\text{ g}$ nước hóa hơi là $1{,}6\text{ phút}$. Bỏ qua hao phí và nhiệt dung của ấm. Tính nhiệt hóa hơi riêng của nước xác định từ thực nghiệm này.
\begin{enumerate}[label=\Alph*.]
  \item $2{,}24 \cdot 10^6\text{ J/kg}$.
  \item $2{,}688 \cdot 10^6\text{ J/kg}$.
  \item $2{,}016 \cdot 10^6\text{ J/kg}$.
  \item $2{,}52 \cdot 10^6\text{ J/kg}$.
\end{enumerate}

\textbf{Câu 17:} Dẫn $20\text{ g}$ hơi nước ở $100^{\circ}\text{C}$ vào bình nhiệt lượng kế chứa $200\text{ g}$ nước ở $20^{\circ}\text{C}$. Bỏ qua nhiệt dung của bình nhiệt lượng kế. Cho $c_{nước} = 4200\text{ J/(kg}\cdot\text{K)}$, $L_{nước} = 2{,}26 \cdot 10^6\text{ J/kg}$. Nhiệt độ cân bằng cuối cùng của hệ là:
\begin{enumerate}[label=\Alph*.]
  \item $60{,}2^{\circ}\text{C}$.
  \item $80{,}5^{\circ}\text{C}$.
  \item $72{,}3^{\circ}\text{C}$.
  \item $78{,}1^{\circ}\text{C}$.
\end{enumerate}

\textbf{Câu 18:} Đổ $1{,}5\text{ lít}$ nước ở nhiệt độ $20^{\circ}\text{C}$ vào một ấm nhôm có khối lượng $600\text{ g}$, đun bằng bếp điện. Sau $35\text{ phút}$ thì có $20\%$ khối lượng nước đã hóa hơi ở $100^{\circ}\text{C}$. Biết rằng chỉ có $75\%$ nhiệt lượng bếp điện tỏa ra được dùng để đun nước và ấm. Cho biết $c_{nước} = 4190\text{ J/(kg}\cdot\text{K)}$, $c_{nhôm} = 880\text{ J/(kg}\cdot\text{K)}$, $L_{nước} = 2{,}26 \cdot 10^6\text{ J/kg}$, khối lượng riêng của nước là $1\text{ kg/lít}$. Công suất cung cấp nhiệt của bếp điện gần giá trị nào nhất?
\begin{enumerate}[label=\Alph*.]
  \item $520\text{ W}$.
  \item $765\text{ W}$.
  \item $620\text{ W}$.
  \item $740\text{ W}$.
\end{enumerate}


\section*{PHẦN II: CÂU HỎI ĐÚNG -- SAI (4 CÂU)}

\textbf{Câu 1:} Cho các phát biểu sau về sự hóa hơi và nhiệt hóa hơi riêng của chất lỏng:
\begin{enumerate}[label=\alph*.]
  \item Nhiệt hóa hơi của một lượng chất lỏng tỉ lệ thuận với khối lượng của phần chất lỏng đã hóa hơi. \hfill (Đ/S)
  \item Nhiệt hóa hơi riêng của một chất lỏng là hằng số và không thay đổi theo nhiệt độ của chất lỏng. \hfill (Đ/S)
  \item Hơi nước ở $100^{\circ}\text{C}$ tỏa ra một lượng nhiệt lớn khi ngưng tụ, nên có thể gây bỏng sâu hơn nước sôi ở cùng nhiệt độ. \hfill (Đ/S)
  \item Đơn vị của nhiệt hóa hơi riêng $L$ là $\text{J/kg}$, trùng với đơn vị của nhiệt nóng chảy riêng $\lambda$. \hfill (Đ/S)
\end{enumerate}

\textbf{Câu 2:} Trong thí nghiệm đo nhiệt hóa hơi riêng của nước bằng nhiệt lượng kế và oát kế:
\begin{enumerate}[label=\alph*.]
  \item Cân điện tử được dùng để xác định khối lượng nước hóa hơi bằng cách cân bình chứa trước và sau đun. \hfill (Đ/S)
  \item Oát kế đo điện năng tiêu thụ và thời gian đun dưới dạng tích số $P \cdot t$. \hfill (Đ/S)
  \item Dây đun điện trở phải được đặt ngập một phần trong nước, không được ngập hoàn toàn để hơi nước thoát ra dễ dàng. \hfill (Đ/S)
  \item Bình nhiệt lượng kế cần có vỏ xốp cách nhiệt để ngăn ngừa nhiệt lượng thất thoát ra môi trường, giảm sai số thực nghiệm. \hfill (Đ/S)
\end{enumerate}

\textbf{Câu 3:} Người ta đun sôi và hóa hơi một phần nước trong chiếc ấm bằng đồng có khối lượng $m_{ấm} = 0{,}4\text{ kg}$ chứa $1{,}5\text{ lít}$ nước ở nhiệt độ ban đầu $20^{\circ}\text{C}$. Sau khi nước sôi được một lúc thì có $0{,}1\text{ lít}$ nước đã biến thành hơi. Cho biết $L_{nước} = 2{,}3 \cdot 10^6\text{ J/kg}$, $c_{nước} = 4200\text{ J/(kg}\cdot\text{K)}$, $c_{đồng} = 380\text{ J/(kg}\cdot\text{K)}$, khối lượng riêng của nước $D = 1\text{ kg/lít}$.
\begin{enumerate}[label=\alph*.]
  \item Nhiệt lượng cần cung cấp để đưa ấm và nước từ $20^{\circ}\text{C}$ lên nhiệt độ sôi $100^{\circ}\text{C}$ là $516{,}16\text{ kJ}$. \hfill (Đ/S)
  \item Nhiệt lượng cần thiết để hóa hơi $0{,}1\text{ lít}$ nước sôi là $230\text{ kJ}$. \hfill (Đ/S)
  \item Nhiệt lượng để đun ấm nhôm/đồng nóng lên lớn hơn nhiệt lượng làm nước nóng lên. \hfill (Đ/S)
  \item Tổng nhiệt lượng đã cung cấp cho ấm nước kể từ ban đầu là $746{,}16\text{ kJ}$. \hfill (Đ/S)
\end{enumerate}

\textbf{Câu 4:} Thực hiện thí nghiệm đưa hơi nước ở $100^{\circ}\text{C}$ ngưng tụ trong một nhiệt lượng kế bằng nhôm khối lượng $150\text{ g}$ chứa $300\text{ g}$ nước ở $15^{\circ}\text{C}$. Nhiệt độ cân bằng cuối cùng đo được là $40^{\circ}\text{C}$. Sau thí nghiệm, khối lượng nước trong bình tăng thêm $12{,}2\text{ g}$. Biết $c_{nước} = 4200\text{ J/(kg}\cdot\text{K)}$, $c_{nhôm} = 900\text{ J/(kg}\cdot\text{K)}$.
\begin{enumerate}[label=\alph*.]
  \item Nhiệt lượng mà nước và bình nhôm thu vào khi nóng lên từ $15^{\circ}\text{C}$ đến $40^{\circ}\text{C}$ là $34\text{ }875\text{ J}$. \hfill (Đ/S)
  \item Khối lượng nước tăng thêm $12{,}2\text{ g}$ chính là khối lượng của hơi nước đã ngưng tụ thành nước. \hfill (Đ/S)
  \item Khi ngưng tụ và giảm nhiệt độ về $40^{\circ}\text{C}$, hơi nước chỉ tỏa ra ẩn nhiệt ngưng tụ mà không tỏa nhiệt khi hạ nhiệt độ. \hfill (Đ/S)
  \item Giá trị nhiệt hóa hơi riêng của nước tính toán được từ thí nghiệm này xấp xỉ $2{,}56 \cdot 10^6\text{ J/kg}$. \hfill (Đ/S)
\end{enumerate}


\section*{PHẦN III: CÂU HỎI TRẢ LỜI NGẮN (6 CÂU)}

\textbf{Câu 1:} Cần cung cấp một nhiệt lượng bằng bao nhiêu MJ để làm hóa hơi hoàn toàn $1{,}5\text{ kg}$ nước ở nhiệt độ sôi $100^{\circ}\text{C}$? (Lấy hằng số $L = 2{,}26 \cdot 10^6\text{ J/kg}$ và làm tròn đến hai chữ số thập phân).
\hrulefill

\textbf{Câu 2:} Đun nóng $400\text{ g}$ cồn từ nhiệt độ $28^{\circ}\text{C}$ đến nhiệt độ sôi $78^{\circ}\text{C}$, sau đó hóa hơi hoàn toàn ở nhiệt độ sôi này. Cho biết nhiệt dung riêng của cồn là $c = 2500\text{ J/(kg}\cdot\text{K)}$ và nhiệt hóa hơi riêng của cồn là $L = 8{,}6 \cdot 10^5\text{ J/kg}$. Tính nhiệt lượng tổng cộng cần cung cấp (theo đơn vị kJ).
\hrulefill

\textbf{Câu 3:} Để xác định nhiệt hóa hơi riêng của cồn, người ta dẫn hơi cồn ở nhiệt độ sôi $T_{sôi} = 78^{\circ}\text{C}$ ngưng tụ trong bình nhiệt lượng kế chứa $250\text{ g}$ nước ở $20^{\circ}\text{C}$. Khối lượng hơi ngưng tụ đo được là $15\text{ g}$. Nhiệt độ cuối cùng của hệ là $32^{\circ}\text{C}$. Bỏ qua nhiệt dung của bình. Tính nhiệt hóa hơi riêng $L$ của cồn theo đơn vị kJ/kg (lấy kết quả nguyên).
\hrulefill

\textbf{Câu 4:} Đun $2\text{ lít}$ nước ở $15^{\circ}\text{C}$ trong ấm bằng nhôm nặng $500\text{ g}$ trên bếp điện có công suất cung cấp nhiệt $1000\text{ W}$. Hiệu suất truyền nhiệt của bếp là $80\%$. Tính thời gian đun (theo phút, làm tròn đến 1 chữ số thập phân) để nước trong ấm sôi và hóa hơi mất $10\%$ khối lượng của nó.
\hrulefill

\textbf{Câu 5:} Tính nhiệt lượng (kJ, làm tròn đến hàng đơn vị) cần cung cấp để $300\text{ g}$ nước đá ở $0^{\circ}\text{C}$ nóng chảy hoàn toàn, sau đó tiếp tục đun sôi đến $100^{\circ}\text{C}$ và hóa hơi $10\%$ khối lượng của nó ở nhiệt độ sôi.
\hrulefill

\textbf{Câu 6:} Trong một thí nghiệm đo $L$ của nước, học sinh ghi lại đồ thị độ giảm khối lượng nước trong bình theo thời gian đun. Từ lúc nước sôi, trong thời gian $5\text{ phút}$ khối lượng nước giảm từ $180\text{ g}$ xuống $150\text{ g}$. Công suất tỏa nhiệt hữu ích của dây đun điện trở đo được là $220\text{ W}$. Tính nhiệt hóa hơi riêng của nước thu được từ thí nghiệm này (theo đơn vị $10^6\text{ J/kg}$, làm tròn đến 2 chữ số thập phân).
\hrulefill

\end{document}